
\newcommand{\bucket}[1]{\textsf{Bucket} #1}

\section{New Lower Bounds for HPC and Multilayer HPC} \label{sec:lower-bound-hpc}

The two main results of this section are an $\Omega(m^2)$ lower-bound against misaligned $r$-round protocols for $\bmhpc_{m,r}$, and an $\Omega(\frac{m^2}{r})$ lower-bound against misaligned $r$-round protocols for $\bhpc_{m,r}$. Notice that in a misaligned $r$-round round protocol for $\bmhpc_{m,r}$ or $\bhpc_{m,r}$, it is players $P_C$ and $P_D$ who begin the protocol by talking with each other. This means that the ``wrong'' pair of players begin to speak, in the sense that they wish to compute the value $\{z_1\} = A_{z_0} \cap B_{z_0}$, but this instance is with Alice and Bob, so they have no way to do this. So the first round of communication says nothing about $z_1$: the best $P_C$ and $P_D$ can do is send some information about all of their Set-Intersection instances, without knowing which is important. So each bit that $P_C$ and $P_D$ communicate with $P_A$ and $P_B$ in the first round, can only reveal $\frac 1 m$ bits of information about the average instance. So, in the next round, $P_A$ and $P_B$ cannot have learned what $C_{z_1}$ or $D_{z_1}$ are, either. But then, how can they learn $z_2$? The difficult situation is now reversed! This ``always one step behind'' situation is similar to what happens for pointer chasing \cite{NisanW93, PonzioRV99, Yehudayoff20}, except now the pointers are ``hidden'' behind set intersection instances.

\medskip\noindent
How much communication must the players use in order to escape this unfortunate situation? In the case of $\bmhpc_{m,r}$, it is not clear how to do this in less than $O(m^2)$ bits of communication. One can think of two strategies: either reveal an entire layer, causing the protocol to become aligned at that layer, or solve roughly $\frac{m}{r}$ set intersections instances in each layer, which will cause the protocol to become aligned (follow the pointer of a solved instance) at any given layer with probability $\frac{1}{r}$. We will prove that these two strategies are essentially optimal.

However, in the case of $\bhpc_{m,r}$, it can be seen that $O(\frac{m^2}{r})$ bits of communication suffice. Indeed, $P_C$ and $P_D$ (who speak first) may choose $N = O(\frac{m}{r})$ of their Set-Intersection instances uniformly at random, solve them all by communicating $O(\frac{m^2}{r})$ bits, and then send the solution of all these instances to $P_A$ and $P_B$. Then, in the next round, $P_A$ and $P_B$ begin the standard protocol for solving $\hpc_{m,r}$, meaning, they compute $z_1$ and send it back, then $P_C$ and $P_D$ compute $z_2$ and send it back, etc. If at any point $P_A$ and $P_B$ compute some $z_j$ which is among the $N$ revealed Set-Intersection instances, then they can skip a step, and learn $z_{j+2}$ immediately without communicating. The protocol is now aligned, and can be finished within the remaining number of rounds. The probability that the path taken crosses one of the instances revealed by $P_C$ and $P_D$ is $\Omega(r/ N)$, which can be made arbitrarily close to $1$ with our choice of $N$.

On the other hand, one should not expect to be able to prove any significant lower-bound when $r$ is a sufficiently large fraction of $\sqrt n$. Indeed, in this case, the birthday paradox says that the path being traversed will cycle with high probability. The players may then follow the path (wasting the first round entirely) until they find a cycle, and then they will immediately learn $z_k$.

\medskip\noindent
The intuitive explanation and the upper-bounds given above will lead us to conjecture the following, necessarily tight lower-bounds, which are the main results of this section.

 \begin{theorem} \label{thm:lb-mhpc}
 Any misaligned $r$-round protocol, solving $\bmhpc_{m,r}$ correctly with constant probability requires $\Omega(m^2)$ bits of communication.
 \end{theorem}
 
 \begin{theorem} \label{thm:lb-hpc}
Any misaligned $r$-round protocol, where $r = o(\sqrt m)$, solving $\bhpc_{m,r}$ correctly with constant probability requires $\Omega(\frac{m^2}{r})$ bits of communication.
 \end{theorem}


 A weaker lower-bound of $\frac{m^2}{r^2}$ for HPC was proven in~\cite{AssadiCK19}.
 Our proofs of these tight lower-bounds follow a similar strategy to the earlier
 proof, but with a somewhat modified protocol inspired by~\cite{Yehudayoff20}.
We begin by proving \Cref{thm:lb-hpc} in \Cref{sec:eps-solving-to-mhpc}, because the proof is slightly simpler, and it is enough for our streaming lower-bounds. The proof of \Cref{thm:lb-hpc} appears in \Cref{sec:eps-solving-to-bhpc}.
%
Both proofs follow a similar outline. We first show that a short, misaligned $r$-round protocol for $\bmhpc$ (or $\bhpc$) must be ``almost solving'' an instance of Set-Intersection on one of the rounds, while revealing little information.
%
The notion of ``almost solving'' will be formally defined in \Cref{sec:almost-solving}. 
%
Then, in \Cref{sec:exact-to-eps}, we show that any low-information protocol for ``almost-solving'' Set-Intersection can be used to actually solve Set-Intersection, also with low information. 
%
Putting it all together in \Cref{sec:put-together}, we will find that a misaligned $r$-round protocol for $\bmhpc_{m,r}$ communicating $o(m^2)$ bits (or for $\bhpc_{m,r}$ communicating $o(\frac{m^2}{r})$ bits), would allow us to solve Set-Intersection with $o(m)$ information cost, which is known to be impossible by \cite{JayramKS03}.


\subsection{``Almost Solving''}\label{sec:almost-solving}

The notion of ``almost solving'' is given precisely below.
Intuitively, a protocol $\eps$-solves Set-Intersection if the distribution of
the intersection point $\Set{\restar} = \rA \cap \rB$ is sufficiently altered by
the execution of the protocol.
Meaning, the distribution of $\restar$ conditioned on knowing the protocol (and
possibly one of the inputs) is sufficiently far-away from the distribution of
$\restar$ without this knowledge.
In~\cite{AssadiCK19}, the chosen notion of \textit{sufficiently far} is the
total variation distance.
Here we will choose positive triangular discrimination.

\begin{definition}[$\eps$-solving Set-Intersection, with respect to positive
  triangular discrimination]\label{def:eps-solving}
Let $\cD$ be a distribution on inputs $(A, B)$ for Set-Intersection.
A randomized protocol $\pi$ \textit{internally $\eps$-solves Set-Intersection} (with
respect to positive triangular discrimination) on $\cD$ if \textit{at least one}
of the following hold true (where $\Pi$ is the transcript of $\pi$, including
public randomness):
\[
  \Exp*_{\rPi, \rA}{\PTD{\mu(\restar \mid \rPi, \rA)}{\mu(\restar \mid \rA)}} \geq \eps
  \qquad\text{ or } \qquad
  \Exp*_{\rPi, \rB}{\PTD{\mu(\restar \mid \rPi, \rB)}{\mu(\restar \mid \rB)}} \geq \eps.
\]
A protocol $\pi$ \textit{externally $\eps$-solves Set-Intersectiom} on $\cD$ if
the following holds true:
\[
  \Exp*_{\rPi}{\PTD{\mu(\restar \mid \rPi)}{\mu(\restar )}} \geq \eps.
\]
\end{definition}

\begin{observation}\label{obs:eps-solving-remove-conditioning}
    From the joint convexity of positive triangular discrimination (\Cref{prop:convex}), it follows that for any (possibly correlated) random variables $\rA, \rB, \rC$,
\begin{align*}
    \E_{\rB,\rC}[\PTD{\mu(\rA\mid \rB,\rC)}{\mu(\rA\mid \rC)}] %
    & = \E_{\rB}\left[\sum_{c} \Pr[\rC =c] \PTD{\mu(\rA\mid \rB,\rC=c)}{\mu(\rA\mid \rC=c)}\right] \\
    & \ge \E_{\rB}\left[\Lambda\left(\sum_{c} \Pr[\rC =c]\mu(\rA\mid \rB,\rC=c), \sum_{c} \Pr[\rC =c]\mu(\rA\mid \rC=c)\right)\right]\\
    & = \E_{\rB}[\PTD{\mu(\rA\mid \rB)}{\mu(\rA)}].
\end{align*}
That is, removing a random variable from both conditions can only decrease the $\Lambda$-distance. In particular, if a protocol externally $\eps$-solves Set-Intersection on some distribution $\cD$, it also internally $\eps$-solves Set-Intersection on $\cD$.
\end{observation}

\begin{observation}\label{obs:eps-solving-private-randomness}
    It may be felt that a more natural definition of internally $\eps$-solving
    should include any eventual private randomness of the protocol $\pi$.
    However, noting that $\{\restar\} = \rA \cap \rB$ is a function of $\rA$ and
    $\rB$, and that $\rA$ and $\rB$ are independent given $\rPi$ (by
    \Cref{fact:rectangle})
    it follows that the distribution of $\restar$ is independent of the private randomness of $\pi$, when given $\rA$ and $\rPi$. So
    \[
    \Exp*_{\rPi, \rA}{\PTD{\mu(\restar \mid \rPi, \rA)}{\mu(\restar \mid \rA)}} = \Exp*_{\rPi, \rA}{\PTD{\mu(\restar \mid \rPi, \rR_A, \rA)}{\mu(\restar \mid \rA)}}
    \]
    where $\rR_A$ is Alice's randomness, and likewise on Bob's side.
    That is, including the private randomness gives the same definition.
\end{observation}


\subsection{Hard Distributions} \label{sec:distribution}


Next, we define two hard distributions for $\cD_\bmhpc$ and $\cD_\bhpc$,  over inputs to $\bmhpc$ and $\bhpc$, respectively. 


\begin{tbox}
    \textbf{Distribution $\cD_\bmhpc$} on tuples $(\rA,\rB,\rC,\rD)$ from the universes $\cX$ and $\cY$: 

    \begin{enumerate}
		\item For any $i\in [r]$, and any $x \in \cX$, sample $(A^i_x,B^i_x) \sim \cD_\SI$ from the universe $\cY$ \emph{independently}. 
		\item For any $i \in [r]$, and any $y \in \cY$, sample $(C^i_y,D^i_y) \sim \cD_\SI$ from the universe $\cX$ \emph{independently}. 
    \end{enumerate}
\end{tbox}

\begin{tbox}
    \textbf{Distribution $\cD_\bhpc$} \cite{AssadiCK19} on tuples $(\rA,\rB,\rC,\rD)$ from the universes $\cX$ and $\cY$: 

    \begin{enumerate}
		\item For any $x \in \cX$, sample $(A_x,B_x) \sim \cD_\SI$ from the universe $\cY$ \emph{independently}. 
		\item For any $y \in \cY$, sample $(C_y,D_y) \sim \cD_\SI$ from the universe $\cX$ \emph{independently}. 
    \end{enumerate}
\end{tbox}

It was shown in~\cite{AssadiCK19} that misaligned randomized $r$-round protocols
for $\hpc_{m,r}$ need $\frac{m^2}{r^2}$ bits of communication under $\cD_\bhpc$.
In \Cref{sec:eps-solving-to-bhpc}, we improve this result and show that
misaligned randomized $r$-round protocols for $\bhpc_{m,r}$ (an easier problem
than $\hpc_{m,r}$) need $\frac{m^2}{r}$ bits of communication.


We restate the following simple observations from~\cite{AssadiCK19} about
$\cD_\bhpc$.
Analogous statements can be made for $\cD_\bmhpc$.

\begin{observation}[\cite{AssadiCK19}]\label{obs:distHPC}
	Distribution $\cD_\bhpc$ is \underline{not} a product distribution. However, in this distribution: 
	\begin{enumerate}[label=(\roman*)]
		\item The inputs to $P_A$ and $P_B$ are {independent} of the inputs to
      $P_C$ and $P_D$, i.e., $(\rA,\rB) \perp (\rC,\rD)$. 
		\item For any $x \in \cX$, $(\rA_x,\rB_x)$ is independent of all other
      $(\rA_{x'},\rB_{x'})$ for $x'\neq x \in \cX$.
      Similarly for all $y,y' \in \cY$ and $(\rC_y,\rD_y)$ and $(\rC_{y'},\rD_{y'})$. 
	\end{enumerate}
\end{observation}

Based on this observation, we also have the following simple property. 

\begin{proposition}[\cite{AssadiCK19}]\label{prop:hpc-rectangle}
  Let $\prot_\bhpc$ be any deterministic protocol for $\bhpc_{m,r}$ on
  $\cD_\bhpc$.
  Then, for any transcript $\Prot$ of $\prot_\bhpc$, $(\rA,\rB)
  \perp (\rC,\rD) \mid \rPi = \Prot$. 
\end{proposition}
\begin{proof}
  Follows from the rectangle property of the protocol $\prot_\bhpc$
  (Fact~\ref{fact:rectangle}).
  In particular, the same exact argument as in the two-player case implies that
  if $[(A, B),(C, D)]$ and $[(A', B'),(C', D')]$ are mapped to the same
  transcript $\Prot$, then $[(A, B),(C', D')]$ and $[(A', B'),(C, D)]$ are
  mapped to $\Prot$ as well.
  Hence, since $(\rA, \rB) \perp (\rC,\rD)$ by \Cref{obs:distHPC}, the
  inputs corresponding to the same protocol would also be independent of each
  other, namely, $(\rA, \rB) \perp (\rC, \rD) \mid \rPi = \Prot$. 
\end{proof}







\subsection{\texorpdfstring{Short Protocols for $\bmhpc$ Give Low-Information $\eps$-Solvers for $\SI$}{Short Protocols for BMHCP Give Low-Information eps-Solvers for Set Intersection}}\label{sec:eps-solving-to-mhpc}

We now show that if we have a short, $r$-round protocol for $\bmhpc_{m,r}$ on $\cD_\bmhpc$, then during one of its rounds it must be externally (and hence also internally) $\eps$-solving Set-Intersection on $\cD_\SI$, while revealing little information. This is proven by making precise the following intuitive line of thinking, the statement of which is formally captured by \Cref{lem:mhpc-dist}. 

\begin{itemize}
  \item
      By the end of round $j$, the speaking pair of players would like to
      have learned $\{\rZ_{j} \} = \rA^j_{\rZ_{j-1}} \cap \rB^j_{\rZ_{j-1}}$
      (for odd $j$, or $\{\rZ_{j} \} = \rC^j_{\rZ_{j-1}} \cap \rD^j_{\rZ_{j-1}}$,
      for even $j$).
      Now suppose that in this round they have very little information about
      $\rZ_{j-1}$.
      Let us say, for example, that the distribution of
      $\rZ_{j-1} \mid \rE_{j-1}, \rPi_{j-1}$ is close to uniform.
      This certainly holds when $j = 2$, since
      $\mu(\rZ_1 | \rE_1, \rPi_1) = \mu(\rZ_1)$ is
      completely uniform (because the wrong pair of players speak first and the
      way we construct $\cD_\bmhpc$).
    


  \item Now suppose that despite this, somehow, the two speaking players
    nonetheless manage to learn a lot of information about $\rZ_{j}$. Then they
    have ``almost solved'' an \textit{unknown} instance (an instance belonging
    to the other players, no less), out of the possible remaining instances of
    Set-Intersection.
    Because they have little information about $\rZ_{j-1}$,
    meaning $\rZ_{j-1}$ is close to uniform, they have thus ``almost solved'' an
    unknown, nearly-uniformly-chosen instance. 

  \item Then, intuitively, they must have ``almost solved'' this unknown instance of Set-Intersection without revealing a lot of information about it. This turns out to be impossible, as we will see in \Cref{sec:exact-to-eps}. 
\end{itemize}


One difficulty in implementing this intuitive strategy is in finding the right measure of information that allows for one to carry through the entire argument. In \cite{AssadiCK19}, \textit{total variation distance from uniform} was used as a measure of how much the speaking players know about which Set-Intersection instance they are trying to solve. But TVD incurs the quadratic loss inherent in Pinsker's inequality, leading to a lower-bound of $\frac{m^2}{r^2}$. The same quadratic loss appeared originally in Nisan and Wigderson's lower-bounds for pointer chasing \cite{NisanW93}, and this led to a three-decade-long gap between known upper and lower-bounds for pointer chasing. This gap was finally closed by Amir Yehudayoff \cite{Yehudayoff20}, using positive triangular discrimination. This suggested the possibility of using Yehudayoff's approach to improve the lower-bound of \cite{AssadiCK19}, and it is this plan that we realize here.

\bigskip\noindent
The following lemma states that no low-information protocol (in the sense of information cost) for $\bmhpc$ can reveal a lot of information (in the sense of $\Lambda$-distance-from-uniform) about a uniformly chosen $\setint$ instance at any layer $j$. The proof shows that any such protocol would give us a low-information protocol for $\eps$-solving $\setint$, which we will show is impossible in \Cref{sec:exact-to-eps}.

\begin{restatable}%[Progress Lemma]
{lemma}{progress-multi} \label{lem:progress}
There is a constant $c > 0$ such that, for any protocol for $\bmhpc_{m,r}$
running on input $(\rA,\rB,\rC,\rD)$ sampled from $\cD_{\bmhpc}$, and for any
$j \in [r]$, we have:
\begin{align*}
\E_{(\rE_j, \rPi_j)}\E_{x \sim \cU_\cX}\left[\Lambda(\mu(\rT^j_x \mid \rE_j,
\rPi_j), \mu(\rT^j_x))\right] 
%
\leq  
%
c \cdot \frac{\mi{\rPi}{\rA^j \mid \rA^{< j}, \rB} + \mi{\rPi}{\rB^j \mid \rB^{> j}, \rA}}{m^2}\\
\E_{(\rE_j, \rPi_j)}\E_{y \sim \cU_\cY}\left[\Lambda(\mu(\rT^j_y \mid \rE_j,
\rPi_j), \mu(\rT^j_y))\right] 
%
\leq  
%
c \cdot \frac{\mi{\rPi}{\rC^j \mid \rC^{< j}, \rD} + \mi{\rPi}{\rD^j \mid \rD^{> j}, \rC}}{m^2}\\
\end{align*}
\end{restatable}

\begin{proof}
    Take any given $j \in [r]$, and consider the following protocol for
    $\eps$-solving $\setint$ on $\cD_{\SI}$:
    
	\begin{abox}
    \textbf{Protocol $\protSI$}: The protocol for $\eps$-solving $\SI$ using a protocol $\protHPC$ for $\bmhpc_{m,r}$. 
    
     \smallskip
    
    \textbf{Input:} An instance $(X, Y) \sim \cD_\SI$ with $X, Y\subseteq\cY$.
    
   % \algline
    
    \begin{enumerate}
    	\item \textbf{Sampling the instance.} Alice and Bob create an instance
        $(A, B, C, D)$ of $\bmhpc_{m,r}$ as follows: %The players create an instance $(A,B)$ of $\SI$ as follows:
	\begin{enumerate}
	
 \item Using \underline{public coins}, Alice and Bob sample an index
   $I \in [m]$ uniformly at random, and  
	Alice sets $A^j_{x_I} = X$ and Bob sets $B^j_{x_I} = Y$ using their given inputs in $\SI$. 
	
 \item Using \underline{public coins}, Alice and Bob sample $A^j_{x_1}, \ldots, A^j_{x_{I-1}}$ from the A-side marginal of $\cD_\SI$, and they sample $B^j_{x_{I+1}}, \ldots, B^j_{x_{m}}$ from the B-side marginal of $\cD_\SI$. 
	
 \item Using \underline{private coins}, Alice samples $A^j_{x_r}$ for $r \in \{I+1, \ldots, n\}$ so that $(A^j_{x_r},B^j_{x_r}) \sim \cD_\SI$. Similarly Bob samples $B^j_{x_\ell}$ for $\ell \in \{1, \ldots, I-1\}$. This completes the construction of $(A^j, B^j)$. 
        
 \item Using \underline{public coins}, Alice and Bob sample $A^{< j}, B^{> j}, C$ and $D$ from their respective marginals in $\cD_\bmhpc$ (this is possible by independence).
 \item Using \underline{private coins}, Alice samples $A^{> j}$ conditioned on the publicly sampled $B^{>j}$, and Bob samples $B^{< j}$ conditioned on the publicly sampled $A^{<j}$. 

        \end{enumerate}
  %      (It should be evident that the distribution of instances $(\bA,\bB,\bC,\bD)$ created by $\protSI$ is $\distHPC$.) 
        
        
	\item The players now run the first $j$ rounds of the protocol $\prot_\bmhpc$ on the instance $(A,B,C,D)$, by Alice playing $P_A$, Bob playing $P_B$, and both Alice and Bob simulating $P_C$ and $P_D$ with no communication (this is possible as both Alice and Bob know $(C, D)$ entirely). 
 
    Let us write $\Pi_\SI$ to denote the messages of $P_A$ and $P_B$, which were actually sent by Alice and Bob, and $\Pi_{\le j}$ to denote the entire simulated transcript of $\prot_\bmhpc$, including the messages of $P_C$ and $P_D$ (which Alice and Bob learned without communicating).
    \end{enumerate}
\end{abox}

From a simple direct-sum argument, we get an upper-bound on the information cost of $\pi_\SI$:

\begin{claim} \label{clm:direct-sum}
    \[ \ICost{\pi_\SI}{\dsi} = \frac{\mi{\rPi_{\le j}}{\rA^j \mid \rA^{< j}, \rB}
    + \mi{\rPi}{\rB^j \mid \rB^{> j}, \rA}}{m}. \]
\end{claim}

\begin{proof}
    The internal information cost of $\prot_\SI$ is
    \[
      \IC{\prot_\SI}{\cD_\SI}
      := \mi{\rX}{\rProtSI \mid \rY, \rR, \rR_B}
      + \mi{\rY}{\rProtSI \mid \rX, \rR, \rR_A},
    \]

    where $R, R_A, R_B$ are the public, Alice's, and Bob's randomness.
    Let us bound the first term, as the second term can be bounded in just the
    same way.
    We then derive:   
\begin{align*}
	\mi{\rPi_\SI}{\rX \mid \rY, \rR, \rR_B} %
    & = \mi{\rPi_{\le j}}{\rX \mid \rY, \rI,\rA^j_{<\rI},\rB^j_{>\rI},\rA^{< j},\rB^{> j},\rC,\rD, \rB^j_{< \rI}, \rB^{< j}}\\
    & = \mi{\rPi_{\le j}}{\rA^j_\rI \mid \rI,\rA^j_{<\rI},\rA^{< j},\rB,\rC,\rD}\\
    & =  \sum_{i=1}^m \frac{1}{m} \mi{\rPi_{\le j}}{\rA^j_i \mid \rI=i,\rA^j_{<i},\rA^{< j},\rB,\rC,\rD}\\
    & =  \sum_{i=1}^m \frac{1}{m} \mi{\rPi_{\le j}}{\rA^j_i \mid \rA^j_{<i},\rA^{< j},\rB,\rC,\rD}\\
    & =  \frac{1}{m} \mi{\rPi_{\le j}}{\rA^j \mid \rA^{< j},\rB,\rC,\rD}\\
    & =  \frac{1}{m} \mi{\rPi_{\le j}}{\rA^j \mid \rA^{< j},\rB}. 
\end{align*}
All equalities are by definition, except for the fourth which is because none of
the remaining distributions depend on the choice of $\rI$, and the last, which
is because $\rC, \rD$ are independent of $\rA^j$, even when conditioned on
$\rPi_{\le j}, \rA^{<j}, \rB$.
\end{proof}

On the other hand, we also get an $\eps$-solver, for $\eps$ being exactly the $\Lambda$ distance-from-uniform of a uniformly chosen coordinate in the $j$-th layer:

\begin{claim} \label{clm:average-coordinate}
$\pi_\SI$ internally $\eps$-solves Set-Intersection on $\dsi$ where:
\[
  \eps \geq \E_{(\rE_j, \rPi_j)}
  \Exp*_{x \sim \cU_\cX} {\PTD{\mu(\rT^j_x \mid \rE_j, \rPi_j)}{\mu(\rT^j_x)}}.
\]
\end{claim}


\begin{proof} 
     By \Cref{def:eps-solving} and \Cref{obs:eps-solving-private-randomness}, we have an $\eps$-solver on Bob's side, for
     \[
       \eps = \E[\Lambda(\mu(\rT \mid \rY, \rPi_\SI,
     \rI, \rR, \rR_B), \mu(\rT \mid \rY, \rI, \rR, \rR_B))], 
     \]
     where $T = X \cap Y = A^j_I \cap B^j_I$,
     $I$ and $R = A^j_{<I}, B^j_{>I}, A^{< j}, B^{> j}, C, D$ are the public
     randomness, and
     $R_B = B^j_{< I}, B^{< j}$ is Bob's private randomness. 
     This is exactly equal to
     \[
     \E[\Lambda(\mu(\rT \mid \rY, \rE_j, \rPi_{j}, \rI, \rR, \rR_B),
     \mu(\rT \mid \rY, \rI, \rR, \rR_B))], \]
     because $E_j, \Pi_j$ is a function of $\Pi_\SI, I, R, R_B$ (Bob knows both
     $A^{<j}$, $B^{<j}$ and $C$, $D$, so he also knows $Z_{< j}$, and
     $\Pi_{\le j}$).
     In turn, this equals
     \begin{align*}
       \E_{\rE_j,\rPi_j, i \sim \uni_{\brac{m}}}
       \Exp*{\PTD{\mu(\rT^j_{x_i} \mid \rE_j, \rPi_{j}, \rI = i, \rY, \rR,\rR_B)}
       {\mu(\rT^j_{x_i}\mid \rY, \rI=i,\rR, \rR_B)}}
     \end{align*}
     Observe that none of the distributions above depend on the event $\rI = i$. So this quantity is equal to
     \begin{align*}
       \E_{\rE_j,\rPi_j, i \sim \uni_{\brac{m}}}
       \Exp*{\PTD{\mu(\rT^j_{x_i} \mid \rE_j, \rPi_{j}, \rY, \rR,\rR_B)}
       {\mu(\rT^j_{x_i}\mid \rY, \rR, \rR_B)}}
     \end{align*}
     Now, by \Cref{obs:eps-solving-remove-conditioning}, this is greater or equal to
     \[
       \E_{\rE_j,\rPi_j, i \sim \uni_{\brac{m}}}
       \Exp*{\PTD{\mu(\rT^j_{x_i} \mid \rE_j, \rPi_{j} )}{\mu(\rT^j_{x_i})}}. \qedhere
     \]
\end{proof}


From \Cref{cor:si-ic-lb-external} proven in the next section, it now follows:
  \begin{align*}
    \frac{\mi{\rPi_{\le j}}{\rA^j \mid \rA^{< j}, \rB} + \mi{\rPi_{\leq j}}{\rB^j \mid \rB^{> j}, \rA}}{m} & \ge c \cdot \eps m\\
    & \geq c m \cdot \E_{(\rE_j, \rPi_j)}\E_{x \sim \cU_\cX}
    \left[\Lambda(\mu(\rT^j_x \mid \rE_j, \rPi_j), \mu(\rT^j_x)) \right].
  \end{align*}
  This concludes the proof of \Cref{lem:progress}.
\end{proof}



\bigskip\noindent
It is now possible to prove that there is no low-information protocol for solving $\bmhpc$. The argument was roughly outlined above: if we assume by induction that the $\Lambda$ distance-from-uniform is small for the answer at the previous layer, then we are essentially trying to solve a uniformly random instance at the current layer. And we just proved that this is not possible.

\begin{lemma}\label{lem:mhpc-dist}
For the same constant $c > 0$ of the \Cref{lem:progress}, and for any protocol
for $\bmhpc_{m,r}$ running on input $(A, B, C, D)$ sampled from
$\cD_{\bmhpc}$, the following statement holds for all $j \in
[r]$:%
\footnote{Intuitively, the statement says that, if $\pi_\bmhpc$
communicates little, then the distribution of $Z_j$, conditioned on the
information available at the end of round $j$, is close to uniform.}
\[
  \E_{(\rE_j, \rPi_j)}\left[\Lambda(\mu(\rZ_j \mid \rE_j, \rPi_j), \mu(\rZ_j))\right]
  \leq 6 c \cdot \frac{\mi{\rPi}{\rA^{\le j},\rC^{\le j} \mid \rB,\rD} + \mi{\rPi}{\rB^{\le j},\rD^{\le j} \mid \rB^{> j}, \rD^{> j}, \rA,\rC}}{m^2}.
\] 

\end{lemma}





This is proven by induction. The base case is easy to see: note that $\mu(\rZ_1
\mid \rE_1, \rPi_1) = \mu(\rZ_1)$ as $\rZ_1$ is independent of $\rE_1 = z_0$
which is fixed and $\rZ_1$ is independent of $\rPi_1$ as $\Pi_1$ is a function
of $(C,D)$.
We obtain case $j$ from case $j - 1$ in two steps: we first show that if the
protocol reveals a lot of information about the variable $\rZ_{j}$, then it must
essentially be $\eps$-solving a \textit{uniformly random instance}. This is
because $Z_{j} = T_{Z_{j-1}}$ and $\rZ_{j-1}$ is, under the inductive case
$j-1$, close to uniform.
We then appeal to the previous lemma.

\begin{proof}
The information---measured by $\Lambda$-distance to uniform---that the $j$-th
round of the protocol reveals about $\rZ_j = \rT_{\rZ_{j-1}}$ is the
information revealed about an instance sampled according to $\rZ_{j-1}$:
  \begin{align*}
		\E_{(\rE_j,\rProt_j)} \Bracket{\Lambda(\mu(\rZ_j \mid
    \rE_j,\rProt_j),\mu(\rZ_j))} &=\E_{(\rE_j,\rProt_j)}
    \Bracket{\Lambda(\mu(\rZ_j \mid \rE_j),\mu(\rZ_j))}\\
		&\hspace{-1.5cm}= \E_{(\rZ_{<j},\rProt_{<j})} \Bracket{\Lambda(\mu(\rZ_j
    \mid \rZ_{<j},\rProt_{<j}),\mu(\rZ_j))} \tag{by definition of
  $\rE_j := (\rZ_{<j},\rProt_{<j})$, and since $\rZ_j$ and $\rPi_j$ are
independent given $\rE_j$.} \\
		&\hspace{-1.5cm}=\E_{(\rZ_{<j},\rProt_{<j})}
    \Bracket{\Lambda(\mu(\rT^j_{\rZ_{j-1}} \mid
    \rZ_{<j-1},\rZ_{j-1},\rProt_{<j}),\mu(\rZ_j))}\tag{by definition, the pointer $Z_j = T_{Z_{j-1}}$}\\
		&\hspace{-1.5cm}= \E_{(\rZ_{<j-1},\rProt_{<j})} \; \E_{z \sim \rZ_{j-1} \mid
    \paren{\rZ_{<j-1},\rProt_{<j}}}\; \Bracket{\Lambda(\mu(\rT^j_{z} \mid
  \rZ_{<j-1},\rProt_{<j}),\mu(\rZ_j))}.
	\end{align*}
  We removed the conditioning on $\rZ_{j-1} = z$ in the first distribution
  because, conditioned on knowing $\rZ_{<j-1},\rProt_{<j}$, for any fixed value of
  $z$ the value of $\rT_{z}$ is independent of the value of $\rZ_{j-1}$:
  \begin{itemize}
	    \item If $j-1$ is odd, $\rT_{z}$ is a function of  $(\rC, \rD)$ and
        $\rZ_{j-1}$ is a function of $(\rA, \rB)$.
            \item If $j-1$ is even, $\rT_{z}$ is a function of $(\rA,\rB)$ and
              $\rZ_{j-1}$ is a function of $(\rC, \rD)$.
            \item And, furthermore, $(\rA,\rB)$ and $(\rB,\rD)$ are independent
              given $\rZ_{<j-1}, \rPi_{<j}$. 
	\end{itemize}

    
	\medskip\noindent
    \Cref{lem:progress} says that if the protocol were to reveal information
    about the solution $\rT_z$ of a truly uniform instance $z$, then this would
    give us a protocol for $\eps$-solving set disjointness. But, by induction,
    we know that the chosen instance $z \sim \rZ_{j-1}$ is close to uniform. It
    should then follow from \Cref{itm:lambda-loss} that if the protocol reveals
    some information about this almost uniform instance, it must also reveal
    information about a truly uniform instance. Let us now make this precise.
    
  Without loss of generality suppose that $j-1$ is odd and hence $z \in \cY$.
  Then we can upper bound the last expression above
  by:\footnote{\label{fn:tvd}In the first inequality it should be understood how
  $\Lambda$ is playing two different roles. On the one hand it is being used as
an information-theoretic quantity (the distance to uniform), and on the other it
is being used in a capacity similar to how total variation distance appears in
the obvious upper-bound $\E_\mu[f(x)] \le \E_\nu[f(x)] + \|\mu - \nu\|_1 \cdot
\max f(x)$.}

	\begin{align*}
          &\E_{(\rE_j,\rPi_j)} \Bracket{\Lambda(\mu(\rZ_j \mid \rE_j,\rPi_j),\mu(\rZ_j))}  \\
          % 
          &\qquad \qquad \leq \E_{(\rZ_{<j-1},\rPi_{<j})} \bracket{6 \cdot \E_{\paren{z \sim \cU_{\cY}}}\Bracket{\Lambda(\mu(\rT^j_{z} \mid \rZ_{<j-1},\rPi_{<j}),\mu(\rZ_j))}} \\
          % 
          &\qquad \qquad  \qquad \qquad \qquad  \qquad + \E_{(\rZ_{<j-1},\rPi_{<j})} \Bracket{\Lambda(\mu(\rZ_{j-1} \mid \rE_{j-1},\rPi_{j-1}),\cU_\cY)} %
            \tag{by \Cref{itm:lambda-loss} and linearity of expectatiom}
          \\
          % 
          &\qquad \qquad = 6 \cdot\E_{(\rZ_{<j-1},\rPi_{<j})} \bracket{ \E_{\paren{z \sim \cU_{\cY}}}\Bracket{\Lambda(\mu(\rT^j_{z} \mid \rZ_{<j-1},\rPi_{<j}),\mu(\rZ_j))}} \\
          &\qquad \qquad  \qquad \qquad \qquad  \qquad + \E_{(\rZ_{<j-1},\rPi_{<j})} \Bracket{\Lambda(\mu(\rZ_{j-1} \mid \rE_{j-1},\rPi_{j-1}),\mu(\rZ_{j-1}))}\tag{by definition $z_{j-1} \in \cY$ and $\mu(\rZ_{j-1}) = \cU_\cY$.}\\
          % 
          &\qquad \qquad\leq \bigg[6 \cdot \underbrace{c \cdot \frac{\mi{\rPi}{\rA^j \mid \rA^{< j}, \rB} + \mi{\rPi}{\rB^j \mid \rB^{> j}, \rA}}{m^2}}_{\text{\Cref{lem:progress}}} \\
          &\qquad \qquad\qquad 
            + \underbrace{6 c \cdot \frac{\mi{\rPi}{\rA^{\le j-1},\rC^{\le j-1} \mid \rB,\rD}
           + \mi{\rPi}{\rB^{\le j-1},\rD^{\le j-1} \mid \rB^{> j-1}, \rD^{> j-1}, \rA,\rC}}{m^2}}_{\text{induction hypothesis (case $j-1$)}}\bigg]\\
    % 
    % 
          & \qquad \qquad= \bigg[6c \cdot \frac{\mi{\rPi}{\rA^j \mid \rA^{< j},\rC^{< j}, \rB,\rD} + \mi{\rPi}{\rB^j \mid \rB^{> j},\rD^{>j}, \rA,\rC}}{m^2} \\
          &\qquad \qquad\qquad 
            + 6 c \cdot \frac{\mi{\rPi}{\rA^{\le j-1},\rC^{\le j-1} \mid \rB,\rD}
            + \mi{\rPi}{\rB^{\le j-1},\rD^{\le j-1} \mid \rB^{> j-1}, \rD^{> j-1}, \rA,\rC}}{m^2}\bigg]\\
          &         \qquad \qquad
          \leq
          \bigg[6 c \cdot \frac{\mi{\rPi}{\rA^j,\rC^j\mid \rA^{< j},\rC^{< j}, \rB,\rD} + \mi{\rPi}{\rB^j,\rD^j \mid \rB^{> j},\rD^{>j}, \rA,\rC}}{m^2} \\
          &\qquad \qquad\qquad 
            + 6 c \cdot \frac{\mi{\rPi}{\rA^{\le j-1},\rC^{\le j-1} \mid \rB,\rD}
            + \mi{\rPi}{\rB^{\le j-1},\rD^{\le j-1} \mid \rB^{> j-1}, \rD^{> j-1}, \rA,\rC}}{m^2}\bigg]\\
          & \qquad\qquad =
            6 c \cdot \frac{\mi{\rPi}{\rA^{\le j},\rC^{\le j} \mid \rB,\rD}
           + \mi{\rPi}{\rB^{\le j},\rD^{\le j} \mid \rB^{> j}, \rD^{> j},
         \rA,\rC}}{m^2},
	\end{align*}
where the last inequality follows from the data processing inequality
(\Cref{fact:it-facts}~\Cref{part:data-processing}).
This concludes the proof.
\end{proof}




\subsection{\texorpdfstring{$\eps$-Solvers for $\SI$ Give Exact Solvers}{eps-Solvers for Set Intersection Give Exact Solvers}}\label{sec:exact-to-eps}



We now show a lower-bound against protocols trying to $\eps$-solve
Set-Intersection, by reducing from exactly solving Set-Intersection:


\begin{theorem}\label{thm:si-ic-lb} For any constant $\gamma \in (0,1)$, there exists a constant $c \ge 1$ such that, given a protocol $\pisi$ which internally 
$\eps$-solves Set-Intersection (in $\ptd$ distance) on $\dsi$ with
$\eps \ge \frac{8}{m}$,
we may obtain a protocol $\pi$ that exactly solves Set-Intersection on $\dsi$,
such that
$\ICost{\pi}{\dsi} \leq {\frac{c}{\eps} \cdot \ICost{\pisi}{\dsi} + \gamma m}$.
\end{theorem}

From the known fact (\Cref{prop:SI-hard}) that Set-Intersection has information
complexity $\Omega(m)$ on $\dsi$ we get the following:

\begin{corollary}\label{cor:si-ic-lb-external}
  There exists a constant $c > 0$ such that
  any protocol $\pisi$ that internally $(\ptd, \eps)$-solves $\setint_m$ on
  inputs from the distribution $\dsi$, with $\eps \geq 8 / m$, has $\ICost{\pisi}{\dsi} \geq c \cdot \eps m$.
\end{corollary}

An analogue of \Cref{thm:si-ic-lb} was proven in
\cite{AssadiR20}, where $\eps$-solving was with respect to total variation
distance instead of triangular discrimination.
It was then shown that any protocol for $\eps$-solving Set-Intersection in this
TVD sense gives a protocol for Set-Intersection with a $O(1/\eps^2)$
blowup in information cost.
\Cref{thm:si-ic-lb}, instead, gives a blowup of $O(1/\eps)$.
So our assumption is stronger, but our conclusion is also stronger.

\begin{proof}[Proof of \Cref{thm:si-ic-lb}]
We begin by recalling the notation we use for the Set-Intersection problem:
Alice and Bob's input sets are
denoted by $X$ and $Y$ respectively, and their corresponding random variables
by $\rX$ and $\rY$.
The unique element in the intersection $X \cap Y$ is denoted by $\estar$,
and its random variable is $\restar$.
Throughout this proof, for a permutation $\sigma$ and a set $S$,
$\sigma(S)$ denotes the set $\Set{ \sigma(s) \mid s \in S }$.

\noindent
Without loss of generality, we may assume that we have a protocol $\pisi$ that
internally $\eps$-solves from \emph{Alice}'s side
Set-Intersection on inputs distributed according to $\dsi$, that is,
\[
  \rPTD = 
  \Exp*_{\rX, \rPi}{
     \PTD*{ \Dist{ \restar \given \rPi, \rX}}
     {\Dist{ \restar, \rX}}
  } \geq \eps,
\]
where $\eps \geq 8 / m$.
We will use $\pisi$ to design a protocol $\pi$ which solves Set-Intersection
on $\dsi$.

For now, suppose that Alice and Bob proceed as follows: they sample a uniformly
random permutation $\Perm$ of $[m]$ using public randomness, and run the
$\eps$-solving protocol $\pisi$ on $(\Perm(X), \Perm(Y))$, obtaining a
transcript $\Pi$.
Alice now knows the induced distribution of the intersection
point $\Dist{\rPerm(\estar) \given \rPi = \Pi, \rPerm(X) = \Perm(X)}$, conditioned
on the transcript she has seen, and Bob knows the analogous distribution with
his input fixed.
Without the conditioning on the transcript, we know by symmetry of $\dsi$ that
$\Dist{\rPerm(\estar) \given \rPerm(X) = \Perm(X)}$ is uniform over $\Perm(X)$.
In fact, we have the following observation about the distribution of inputs
generated by applying $\rPerm$.

\begin{observation}\label{obsv:dsi}
  For any fixed $X$ and $Y$ such that $\card{X} = \card{Y} = m / 4$, and
  $\card{X \cap Y} = 1$,
  and a uniformly random permutation $\rPerm$,
  the instance $(\rPerm(X), \rPerm(Y))$ is distributed according to $\dsi$.
\end{observation}

To see this, consider the fixed values $X = \Set{1, \ldots m / 4}$ and
$Y = \Set{m / 4, \ldots m / 2 - 1}$.
Because $\rPerm$ is a uniformly random permutation, it maps $m / 4$ to a
uniformly random element of $\brac{m}$, and the rest of $X$ and $Y$ to
uniformly random subsets of size $m / 4 - 1$.
By symmetry, this argument applies to any fixed $X$ and $Y$ of size $m / 4$
intersecting in exactly one element.


\subsubsection{A single round of \texorpdfstring{$\pi$}{the protocol}}\label{subsec:single-round}

We start by describing a single round of the protocol $\pi$.
On the input $(X, Y)$:
\begin{enumerate}
  \item Alice and Bob publicly sample a random permutation $\Perm$ of
    $\bracket{m}$, and apply it to $X$ and $Y$ to obtain $\Perm(X)$ and
    $\Perm(Y)$ respectively.
  \item They run $\pisi$ on $\Perm(X)$ and $\Perm(Y)$ to obtain a transcript
    $\Pi$.
    By doing so, Alice learns a distribution $\qpi_1, \ldots , \qpi_m$, where
    \[ \qpi_i = \Prob{ \rPerm(\estar) = i \given \rPi = \Pi, \rPerm(X) = \Perm(X)} . \]
  \item Alice uses $\qpi$ to assign ``scores'' to each element $e$ of
    $\brac{m}$, which reflect the likelihood that $\Set{e}$ is the
    intersection of $X$ and $Y$.
\end{enumerate}

This strategy follows the general idea of~\cite{AssadiR20}, but where they assigned the score $1$ to all $e$ such that $\qpi_{\Perm(e)}$ is in the largest-probability half of $\qpi_1, \ldots , \qpi_m$ (and this works for $\eps$-solving under TVD), we assign scores in a smoother fashion.
Let $i = \Perm(e)$ and $n = m / 4$ for brevity.
Then if $\qpi_{i} \leq 1/n$, we give $e$ a score of $0$.
Otherwise, its score is $(\qpi_i - 1/n) / (\qpi_i + 1/n)$.
This choice of the ``scoring function'' may seem odd at this juncture, so
we offer two comments about it:
\begin{itemize}
  \item Note that this is equal to $1 - (2 / n) / (\qpi_i + 1/n)$, which is an
    increasing function of $\qpi_i$, and hence passes a basic sanity test
    (i.e.\ if $\pisi$ says an element is more likely to be $\Perm(\estar)$,
    we believe it).
  \item 
    The scoring function was chosen to make the expected score of $\estar$
    larger than the expected score of any non-$\estar$ element by at least
    $\rPTD$, and the variance of every element's score turned out to be small.
    %It is a mystery to us why the math works out.
    %\parth{I want to leave this last line in.}
\end{itemize}

We package the properties of this scoring function into a few claims below.
Throughout these claims, we will use the shorthand notation
$\rA \mid B$ to mean $\rA \mid \rB = B$, where $\rA$ and $\rB$ are random
variables, and $B$ is some fixed value that $\rB$ can take.

\begin{claim}\label{clm:exp-score-diff}
  For any non-$\estar$ element $e \in \brac{m}$,
  \[ \Exp_{\rPerm} {\score(\estar) - \score(e)} \geq \rPTD / 2. \]
\end{claim}

\begin{proof}
  The proof is in two parts; first we bound $\score(\estar)$ under some
  appropriate conditioning, and then $\score(e)$ for non-$\estar$ elements.

  \noindent
  Fix any $(X,Y)$ with $|X| = |Y| = n$ and $|X\cap Y| = 1$.
  Consider a fixed choice of $\Perm(X)$ and the transcript $\Pi$, and hence
  a distribution $\qpi$.
  Let $P$ denote the set of indices where $\qpi_i > 1/n$ (i.e.\ those that lead
  to positive scores).
  Note that $P \subseteq \Perm(X)$; because we condition on $\Perm(X)$
  (and more) to obtain $\qpi$.
  By definition, the probability that $\estar$ is mapped to $i$ is $q_i$, and
  hence:
  \[
    \Exp*_{\rPerm \given \Perm(X), \Pi} {\score(\estar)}
    = \sum_{i \in P} \qpi_i \cdot \frac{\qpi_i - 1/n}{\qpi_i + 1/n}.
  \]
  On the other hand, fixing both $\Perm(X)$, $\Perm(Y)$ (and hence
  also $\Perm(\estar)$) and $\Pi$, fixes the distribution $\qpi$,
  but leaves completely free the map
  $\Perm: X \setminus \Set{\estar} \to \Perm(X) \setminus \Set{\Perm(\estar)}$.
  So we find that:
 
    $\rPerm(e)$ is uniformly random over
      $\Perm(X) \setminus \Set{\Perm(\estar)}$ for any non-$\estar$ element
      $e \in X$.
  
  
  This means that for any $i \in \brac{m}$,
  \[
    \Prob_{ \rPerm \given \Perm(X), \Perm(Y), \Pi } {\rPerm(e) = i} \leq
    \frac 1{n - 1}.
  \]
  Hence we can compute the expected score of $e$ as:
  \begin{align*}
    \Exp_{\rPerm \given \Perm(X), \Perm(Y), \Pi}{\score(e)}
    &= \sum_{i \in P}  \Prob{\rPerm(e) = i} \cdot \frac{\qpi_i - 1/n}{\qpi_i + 1/n}
    \leq \sum_{i \in P}  \frac 1{n - 1} \cdot \frac{\qpi_i - 1/n}{\qpi_i + 1/n}\\
    &= \sum_{i \in P}  \pparen*{\frac 1{n} + \frac 1{n(n - 1)}} \cdot \frac{\qpi_i - 1/n}{\qpi_i + 1/n}
    \leq \frac 1n + \sum_{i \in P}  \frac 1{n} \cdot \frac{\qpi_i - 1/n}{\qpi_i + 1/n},
  \end{align*}
  where the last inequality follows because $\card{P} \leq n - 1$.
  And now taking expectation over $\Perm(X)$ and $\Pi$:
  \begin{align*}
    \expect_{\rPerm(X), \rPi}
    \Exp_{\rPerm \given \rPi, \rPerm(X)}{\score(\estar) - \score(e)}
    &=
    \Exp*_{\rPerm(X), \rPi}{
    \Exp_{\rPerm \given \rPi, \rPerm(X)}{\score(\estar)}
    - \expect_{\rPerm(Y) \given \rPi, \rPerm(X)} \Exp_{\rPerm \given \rPerm(X), \rPerm(Y)}{\score(e)}
    }\\
    &\geq
    \Exp*_{\rPerm(X), \rPi}{\sum_{i \in P} \frac{ \paren{\qpi_i - 1/n}^2 }{\qpi_i + 1/n} - 1/n} \\
    &=
    \Exp*_{\rPerm(X), \rPi}{\PTD{\qpi}{\uni_{\Perm(X)}}} - 1/n
    = \rPTD - 1 / n \geq \rPTD / 2.
    \intertext{where the last inequality is because
      $\rPTD \geq \eps \geq 8 / m = 2 / n$.
    }
  \end{align*}
\end{proof}

\begin{claim}\label{clm:variance-score-target}
  The variance of $\score(\estar)$ is at most $\rPTD$.
\end{claim}

\begin{proof}
  Recall that $\Var{X} \leq \Exp{X^2}$ for any random variable $X$.
  As before, we will bound $\Exp{\score(\estar)}$ by conditioning on
  fixed choices of $\Perm(X)$ and the transcript $\Pi$, which
  gives Alice the distribution $\qpi$.
  Also as before, let $P$ denote the set of indices $i$ such that
  $q_i \geq 1 / n$.
  \begin{align*}
    \Exp*_{\rPerm \mid \Perm(X), \Pi}{{\score(\estar)}^2}
    &= \sum_{i \in P} \qpi_i \cdot \pparen*{\frac {\qpi_i - 1 / n}{\qpi_i + 1 / n}}^2\\
    &= \sum_{i \in P} \frac {\qpi_i}{\qpi_i + 1 / n} \cdot \frac{\pparen{\qpi_i - 1/n}^2}{\qpi_i + 1 / n}\\
    &\leq \PTD{\qpi}{ \uni_{\Perm(X)}}.
  \end{align*}
  And taking expectation over $\Perm(X)$ and $\Pi$:
  \[
    \expect_{\rPerm(X), \rPi}
    \Exp*_{\rPerm \mid \rPerm(X), \rPi}{{\score(\estar)}^2}
    \leq \Exp*_{\rPerm(X), \rPi}{\PTD{\qpi}{\uni_{\rPerm(X)}}}
    = \rPTD.
  \]
\end{proof}

\begin{claim}\label{clm:variance-score-other}
  For any $e (\neq \estar) \in X$, the variance of $\score(e)$ is at most
  $2\rPTD$.
\end{claim}

\begin{proof}
  Again, because $\Var{X} \leq \Exp{X^2}$ for any random variable $X$, and
  conditioning on fixed choices of $\Perm(X), \Perm(Y)$ and $\Pi$,
  $\rPerm(e)$ for $e \in X \setminus \Set{\estar}$ is uniformly random over
  $\Perm(X \setminus \Set{\estar})$:
  \begin{align*}
    \Exp*_{\rPerm \mid \Perm(X), \Perm(Y), \Pi}{{\score(e)}^2}
    &= \sum_{i \in P} \frac 1{n - 1}
    \cdot \pparen*{\frac {\qpi_i - 1 / n}{\qpi_i + 1 / n}}^2\\
    &\leq \sum_{i \in P} \frac 1{n - 1}
    \cdot \frac{\paren{\qpi_i - 1/n}^2}{\qpi_i + 1 / n} \cdot \frac 1{1 / n}\\
    &= \PTD*{\qpi}{\uni_{\Perm(X)}} \cdot \pparen*{1 + \frac 1{n - 1}}
    \leq 2\PTD*{\qpi}{ \uni_{\Perm(X)}},
  \end{align*}
  where the first inequality holds because $\qpi_i \geq 0$.
  And taking expectation over $\Perm(X), \Perm(Y)$ and $\Pi$:
  \[
    \Exp*_{\rPerm(X), \rPi}{
      \Exp*_{\rPerm(Y) \mid \rPerm(X), \rPi} {
        \Exp*_{\rPerm \mid \rPerm(X), \rPerm(Y), \rPi}{{\score(e)}^2}
      }
    }
    \leq
    \Exp*_{\rPerm(X), \rPi}{
      2\PTD*{\qpi}{\uni_{\rPerm(X)}}
    }
    = 2\rPTD.
  \]
\end{proof}

\subsubsection{The actual protocol}\label{subsec:actual-protocol}

We are ready to describe the protocol $\pi$.
Let $\gamma \in (0,1)$ be any constant.
We will run $k = 1600 / (\eps \gamma^2)$ (independent) rounds of the
protocol in \Cref{subsec:single-round}, and for each element $e \in X$,
Alice will add up its scores in each round to obtain the total score $\totscore(e)$. \Cref{clm:exp-score-diff} tells us that we should expect the total score for $\estar$ to be higher than the total score of an element $e \neq \estar$.
We then choose the threshold
\[
  \tau = \paren{\Exp{\totscore(\estar)} + \Exp{\totscore(e)}} / 2,
\]
to be exactly in the middle of these two expected values.

Alice will then send to Bob all elements in $X$ which have $\totscore$ at least
$\tau$; let us call this set $S$.
Bob will then compute the intersection $S \cap Y$, and report the unique element
that is (hopefully) present there as the answer.

To complete the reduction, we need to show that:
\begin{itemize}
  \item The total score of the target element $\totscore(\estar)$ is at least
    $\tau$.
    This will ensure that $\estar \in S$, which means that $\pi$ is correct.
  \item The number of elements Alice sends is \emph{small}.
    This is required to control the information cost of $\pi$.
\end{itemize}
We will get both of these properties only with some (constant) high probability;
in particular, if $S$ is too large, Alice will just report failure instead
of communicating $S$.

\begin{claim}
  $\totscore(\estar) \leq \tau$ with probability at most $1 / 100$.
\end{claim}

\begin{proof}
Using Chebyshev's inequality on the total score of $\estar$, we have:

\begin{align*}
  &\mkern-36mu\Prob*{ \card*{\totscore(\estar) - \Exp{\totscore(\estar)}}
  > 1/2 \cdot \Exp{\totscore(\estar) - \totscore(e)} } \\
  &\leq \frac{ \Var*{\totscore(\estar)} }{\pparen*{1/2 \cdot \Exp{\totscore(\estar) - \totscore(e)}}^2}
  = \frac{k \cdot \Var{\score(\estar)}}{k^2 / 4 \cdot \Exp{\score(\estar) - \score(e)}^2}\\
  &\leq \frac{4\rPTD}{k\cdot\pparen{\rPTD/2}^2}
  \leq 16 / k\eps = \gamma^2/100,
\end{align*}
where the second inequality follows from \Cref{clm:exp-score-diff,clm:variance-score-target}.
Hence the probability that $\totscore(\estar)$ is less than $\tau$ is at most
$1/ 100$.
\end{proof}

\begin{claim}
  For any non-$\estar$ element $e \in X$, $\totscore(e) \geq \tau$ with
  probability at most $\gamma^2 / 50$.
\end{claim}

\begin{proof}
Using Chebyshev on the total score of $e \neq \estar$, we have:

\begin{align*}
  &\mkern-36mu\Prob*{ \card*{\totscore(e) - \Exp{\totscore(e)}}
  > 1/2 \cdot \Exp{\totscore(\estar) - \totscore(e)} } \\
  &\leq \frac{ \Var*{\totscore(e)} }{\pparen*{1/2 \cdot \Exp{\totscore(\estar) - \totscore(e)}}^2}
  = \frac{k \cdot \Var{\score(e)}}{k^2 / 4 \cdot \Exp{\score(\estar) - \score(e)}^2}\\
  &\leq \frac{4 \cdot 2\rPTD}{k \cdot \paren{\rPTD / 2}^2}
  \leq 32 / k\eps = \gamma^2/50,
\end{align*}
where the second inequality follows from \Cref{clm:exp-score-diff,clm:variance-score-other}.
\end{proof}

This means that the expected size of $S \setminus \Set{\estar}$ is at most
$\gamma^2 m / 50$, and hence by Markov's this number is $> \gamma^2 m / 10$
with probability $\leq 1/5$.
If the number of elements above the threshold exceeds $\gamma^2 m / 10 + 1$,
then the protocol $\pi$ fails.

And now we have everything we need to bound the information cost of $\pi$.
Note that this part of the argument is exactly as in~\cite{AssadiR20}.
By the chain rule of mutual information, the information cost of $\pi$
on $\dsi$ can be broken down into three components:
\begin{itemize}
  \item The mutual information
    $\Inf{\Perm_1, \ldots , \Perm_k, R_1, \ldots , R_k ; \rX \given \rY}$
    and the symmetric term with $\rX$ and $\rY$ exchanging roles,
    where $\Perm_i$ denotes the random permutations used to scramble the input
    in the $i-$th round of $\pi$, and $R_i$ denotes the randomness of the
    $i$-th run of $\pisi$.
    These are both $0$, because the permutations and the randomness are chosen
    independently of the inputs $\rX$ and $\rY$.
  \item The mutual information
    $\Inf{\rPi_1, \ldots , \rPi_k ; \rX \given \rY, \Sigma, \rR}$
    (and the symmetric term with $\rX$ and $\rY$ swapped), where
    $\rPi_1, \ldots , \rPi_k$ denote the transcripts of the $k$ runs
    of $\pisi$, $\Sigma = \Perm_1, \ldots , \Perm_k$, and
    $\rR = R_1, \ldots , R_k$.
    By the chain rule of mutual information, we have:
    \[ \Inf{\rPi_1, \ldots , \rPi_k ; \rX \given \rY, \Sigma, \rR}
      \leq \sum_{i = 1}^k \Inf{\rPi_i ; \rX \given \rY, \Sigma, \rR, \rPi_{< i}}.
    \]
    Observe that conditioning on $\rX$, $\rY$, $\Perm_i$, and $R_i$ fixes
    $\rPi_i$, and hence makes it independent of $\rPi_{< i}$, $\Perm_{-i}$,
    and $R_{-i}$.
    Applying \Cref{prop:info-decrease}, we get that the $i$-th term of the
    sum above is upper bounded by $\Inf{\rPi_i ; \rX \given \rY, \Perm_i, R_i}$.
    But now (once again) $\rPi_i$ is fixed after conditioning on $\Perm_i(\rX)$,
    $\Perm_i(\rY)$, and $R_i$, and hence independent of the remaining parts
    of $\Perm_i$ and $\rX$ and $\rY$.
    So by \Cref{prop:info-decrease}, each term in the sum is bounded above by
    $\Inf{\rPi_i ; \Perm_i(\rX) \given \Perm(\rY), R_i}$.
    Repeating this entire argument for the symmetric term tells us that:
    \[ \Inf{\rPi_1, \ldots , \rPi_k ; \rX \given \rY, \Sigma, \rR} +
      \Inf{\rPi_1, \ldots , \rPi_k ; \rY \given \rX, \Sigma, \rR}
      \leq \sum_{i = 1}^k
      \Inf{\rPi_i ; \Perm_i(\rX) \given \Perm(\rY), R_i} +
      \Inf{\rPi_i ; \Perm_i(\rY) \given \Perm(\rX), R_i},
    \]
    which is exactly $k \cdot \IntIC_{\dsi}(\pisi)$.
  \item The mutual information between the final set $\rS$ Alice sends and
    her input.
    This is trivially upper bounded by $\Ent{\rS}$, which (since $\rS$ is a
    set of size at most $\gamma^2 \cdot m / 10$) is upper bounded by
    $H_2(\gamma^2 / 10) \cdot m$.
    Since $H_2(p) \leq 2\sqrt{p \cdot (1 - p)} \leq 2\sqrt{p}$
    (by \Cref{fact:ent-ub}),
    this is at most $\gamma m$.
\end{itemize}
So letting $c = \frac{1600}{\gamma^2}$, we get that $\IntIC_{\dsi}(\pi) \le \frac{c}{\eps} \IntIC_{\dsi}(\pisi) + \gamma m$.
\end{proof}

\begin{remark}
  The protocol $\pi$ actually works in the worst case, that is, it needs no
  assumptions about $X$ and $Y$ other than $\card{X} = \card{Y} = m / 4$, and
  that $\card{X \cap Y} = 1$.
  But of course without a distribution on its inputs we cannot talk about
  its information cost.
\end{remark}





\subsection{\texorpdfstring{Short Protocols for $\bhpc$ Give Low-Information $\eps$-Solvers for $\SI$}{Short Protocols for BHCP Give Low-Information eps-Solvers for Set Intersection}}\label{sec:eps-solving-to-bhpc}



We also wish to prove a tight lower bound for $\bhpc_{m,r}$. Note that this lower-bound is not necessary to derive our lower-bounds for $k$-cores and degeneracy, so this section may be skipped for that purpose. The lower-bound ultimately follows from the following lemma, which is analogous to \Cref{lem:mhpc-dist}.

\begin{lemma}\label{lem:nonmhpc-dist}
There is a constant $c > 0$ such that the following statement holds for all $j
\in [r]$:\footnote{Intuitively, the statement says that, if $\pi_\bhpc$
communicates little, then the distribution of $\rZ_j$, conditioned on the
information available at the end of round $j$, is close to uniform.}
\[
\mathsf{Hypothesis}_j:\;\;\;\;\;\;\;\;\E_{(\rE_j, \rPi_j)}\left[\Lambda(\mu(\rZ_j \mid \rE_j, \rPi_j), \mu(\rZ_j))\right] \leq j \cdot c \cdot \left( \frac{CC(\pi_\bhpc)+r\log m}{m^2} + \frac{r} m\right).
\] 

\end{lemma}



One can already see that the statement is more complicated. This is owed to the possibility of the pointers pointing either to some of the (few) instances about which the protocol has revealed information, and the possibility of the pointers looping. These two obstacles raise technical difficulties, which are overcome in two ways: the number of jumps $r$ is assumed to be smaller than $\sqrt m$, so looping is unlikely, and there are only $O(m/r)$ instances about which the protocol could have revealed significant information, hence when hiding an instance of set-intersection in a random position, as we did before, it is unlikely we will hit such an instance.

As before, this is proven by induction. The base case $\mathsf{Hypothesis}_1$ is easy to see: note that $\mu(\rZ_1 \mid \rE_1, \rPi_1) = \mu(\rZ_1)$ as $\rZ_1$ is independent of $\rE_1 = z_0$ which is fixed and $\rZ_1$ is independent of $\rPi_1$ as $\Pi_1$ is a function if $(C,D)$. We obtain
$\mathsf{Hypothesis}_j$ from $\mathsf{Hypothesis}_{j-1}$ in two steps: we first show that if the protocol reveals a lot of information about the variable $\rZ_{j}$, then it must essentially be ``almost-solving'' a \textit{uniformly random instance}. This is because $Z_{j} = T_{Z_{j-1}}$ and $\rZ_{j-1}$ is, under $\mathsf{Hypothesis}_{j-1}$, close to uniform. We then show that this is impossible for a short protocol to do, namely: 

%\begin{lemma}[Progress Lemma] \label{lem:progress}
\begin{restatable}%[Progress Lemma]
{lemma}{nonmprogress} \label{lem:nonmprogress}
There is a constant $c' > 0$ such that for any $j \in [r]$, we have:
\begin{align*}
\E_{(\rE_j, \rPi_j)}\E_{x \sim \cU_\cX}\left[\Lambda(\mu(\rT_x \mid \rE_j, \rPi_j), \mu(\rT_x))\right] \leq  c' \cdot \left( \frac{CC(\pi_\bhpc)+r\log m}{m^2} + \frac{j} m\right)\\
\E_{(\rE_j, \rPi_j)}\E_{y \sim \cU_\cY}\left[\Lambda(\mu(\rT_y \mid \rE_j, \rPi_j), \mu(\rT_y))\right] \leq c' \cdot \left( \frac{CC(\pi_\bhpc)+r\log m}{m^2} + \frac{j} m\right).
\end{align*}
\end{restatable}

We state the proof of \Cref{lem:nonmprogress} shortly. Before that, we show how \Cref{lem:nonmprogress} lets us complete the induction in order to prove \Cref{lem:nonmhpc-dist}. This also follows very closely to the analogous proof of \cite{AssadiCK19} where they showed it for total variation distance.

\begin{proof}[Proof of \Cref{lem:nonmhpc-dist}]
The information---measured by $\Lambda$-distance to uniform---that the $j$-th round of the protocol reveals about $\rZ_j = \rT_{\rZ_{j-1}}$ is the information revealed about an instance sampled according to $\rZ_{j-1}$:
  \begin{align*}
		\E_{(\rE_j,\rProt_j)} \Bracket{\Lambda(\mu(\rZ_j \mid \rE_j,\rProt_j),\mu(\rZ_j))} &=\E_{(\rE_j,\rProt_j)} \Bracket{\Lambda(\mu(\rZ_j \mid \rE_j),\mu(\rZ_j))}\\
		&\hspace{-1.5cm}= \E_{(\rZ^{<j},\rProt^{<j})} \Bracket{\Lambda(\mu(\rZ_j \mid \rZ^{<j},\rProt^{<j}),\mu(\rZ_j))} \tag{by definition of $\rE_j := (\rZ^{<j},\rProt^{<j})$, and since $\rZ_j$ and $\Pi_j$ are independent given $\rE_j$.} \\
		&\hspace{-1.5cm}=\E_{(\rZ^{<j},\rProt^{<j})} \Bracket{\Lambda(\mu(\rT_{\rZ_{j-1}} \mid \rZ^{<j-1},\rZ_{j-1},\rProt^{<j}),\mu(\rZ_j))}\tag{by definition, the pointer $Z_j = T_{Z_{j-1}}$}\\
		&\hspace{-1.5cm}= \E_{(\rZ^{<j-1},\rProt^{<j})} \; \E_{z \sim \rZ_{j-1} \mid \paren{\rZ^{<j-1},\rProt^{<j}}}\; \Bracket{\Lambda(\mu(\rT_{z} \mid \rZ^{<j-1},\rProt^{<j}),\mu(\rZ_j))}.
	\end{align*}
	We removed the conditioning on $\rZ_{j-1}$ in the first distribution because, conditioned on knowing $\rZ^{<j-1},\Prot^{<j}$, for any fixed value of $z$ the value of $\rT_{z}$ is independent of the value of $\rZ_{j-1}$: \begin{itemize}
	    \item If $j-1$ is odd, $\rT_{z}$ is a function of  $(\rC,\rD)$ and $\rZ_{j-1}$ is a function of $(\rA,\rB)$.
            \item If $j-1$ is even, $\rT_{z}$ is a function of $(\rA,\rB)$ and $\rZ_{j-1}$ is a function of $(\bC,\bD)$.
            \item And, furthermore, $(\rA,\rB)$ and $(\rB,\rD)$ are independent given $\rZ^{<j-1}, \rPi^{<j}$. 
	\end{itemize}

    
	\medskip\noindent
    \Cref{lem:nonmprogress} says that if the protocol were to reveal information about the solution $\rT_z$ of a truly uniform instance $z$, then this would give us a protocol for $\eps$-solving set disjointness. But, by induction, we know that the chosen instance $z \sim \rZ_{j-1}$ is close to uniform. It should then follow from \Cref{itm:lambda-loss} that if the protocol reveals some information about this almost uniform instance, it must also reveal information about a truly uniform instance. Let us now make this precise.
    
	Without loss of generality suppose that $j-1$ is odd and hence $z \in \cY$. Then 
	we can upper bound the last expression above by:
	\begin{align*}
		&\E_{(\rE_j,\rProt_j)} \Bracket{\Lambda(\mu(\rZ_j \mid \rE_j,\rProt_j),\mu(\rZ_j))}  \\
  %
		&\qquad \qquad \leq \E_{(\rZ^{<j-1},\rProt^{<j})} \bracket{6 \cdot \E_{\paren{z \sim \cU_{\cY}}}\Bracket{\Lambda(\mu(\rT_{z} \mid \rZ^{<j-1},\rProt^{<j}),\mu(\rZ_j))}} \\
  %
		 &\qquad \qquad  \qquad \qquad \qquad  \qquad + \E_{(\rZ^{<j-1},\rProt^{<j})} \Bracket{\Lambda(\mu(\rZ_{j-1} \mid \rE_{j-1},\rProt_{j-1}),\cU_\cY)} %
   \tag{by \Cref{itm:lambda-loss} and linearity of expectatiom}
   \\
   %
   &\qquad \qquad = 6 \cdot\E_{(\rZ^{<j-1},\rProt^{<j})} \bracket{ \E_{\paren{z \sim \cU_{\cY}}}\Bracket{\Lambda(\mu(\rT_{z} \mid \rZ^{<j-1},\rProt^{<j}),\mu(\rZ_j))}} \tag{$\ast$}\label{eq:final-inequality}\\
		 &\qquad \qquad  \qquad \qquad \qquad  \qquad + \E_{(\rZ^{<j-1},\rProt^{<j})} \Bracket{\Lambda(\mu(\rZ_{j-1} \mid \rE_{j-1},\rProt_{j-1}),\mu(\rZ_{j-1}))}\tag{by definition $z_{j-1} \in \cY$ and $\mu(\rZ_{j-1}) = \cU_\cY$.}\\
   %
   &\qquad \qquad\leq 6 \cdot \underbrace{c' \cdot \left( \frac{CC(\pi_\bhpc)+r\log m}{m^2} + \frac{j} m\right)}_{\text{\Cref{lem:nonmprogress}}} + \underbrace{(j-1) \cdot c \cdot \left( \frac{CC(\pi_\bhpc)+r\log m}{m^2} + \frac{r} m\right)}_{\mathsf{Hypothesis}_{j-1}}\\
		%
  & \qquad \qquad\leq j \cdot c \cdot \left( \frac{CC(\pi_\bhpc)+r\log m}{m^2} + \frac{r} m\right) \\
  &\tag{where we replaced $j \leq r$ by $r$ in the first term and set $c = 6c'$}
	\end{align*}

This concludes the proof.
\end{proof}


\bigskip\noindent
We now prove \Cref{lem:nonmprogress}. The proof shows that any protocol that reveals information about the solution $T_z$ to random instance $z$ can be converted to a low-information protocol for externally $\eps$-solving Set-Intersection. This will later be shown to be impossible. Let us state the lemma again:

\nonmprogress*
Suppose towards a contradiction that this equation does not hold. We use $\protHPC$ to design a protocol $\protSI$
	that can externally $\eps$-solve the Set-Intersection problem $(A_x,B_x)$ for a uniformly at random chosen $x \in \cX$ and appropriately chosen $\eps \in (0,1)$ to be determined later (see \Cref{def:eps-solving} for the notion of $\eps$-solving). 
	
	\begin{tbox}
    \textbf{Protocol $\protSI$}: The protocol for $\eps$-solving $\SI$ using a protocol $\protHPC$ for $\bhpc_{m,r}$. 
    
     \smallskip
    
    \textbf{Input:} An instance $(X, Y) \sim \dsi$ over the universe $\cY$ and
    a number $j \in [r]$.
    
   % \algline
    
    \begin{enumerate}
    	\item \textbf{Sampling the instance.} Alice and Bob create an instance
        $(A, B, C, D)$ of $\bhpc_{m,r}$ as follows: %The players create an instance $(A,B)$ of $\SI$ as follows:
	\begin{enumerate}
	\item Using \underline{public coins}, Alice and Bob sample an index $i \in [m]$ uniformly at random, and  
	Alice sets $A_{x_i} = X$ and Bob sets $B_{x_i} = Y$ using their given inputs in $\SI$. 
	\item Using \underline{public coins}, Alice and Bob sample $A_{x_1}, \ldots, A_{x_{i-1}}$ from the A-side marginal of $\cD_\SI$, and they sample $B_{x_{i+1}}, \ldots, B_{x_{m}}$ from the B-side marginal of $\cD_\SI$. 
	\item Using \underline{private coins}, Alice samples $A_{x_r}$ for $r \in \{i+1, \ldots, n\}$ so that $(A_{x_r},B_{x_r}) \sim \cD_\SI$. Similarly Bob samples $B_{x_\ell}$ for $\ell \in \{1, \ldots, i-1\}$. This completes construction of $(\bA,\bB)$. 
  \item Using \underline{public coins}, Alice and Bob sample $(C, D)$ completely
    from $\cD_\bhpc$ (this is possible by Observation~\ref{obs:distHPC} as
    $(\rA, \rB) \perp (\rC, \rD)$). 
        \end{enumerate}
  
        
           \item\label{line:AB} \textbf{Computing the answer.} 
          For $q = 1$ to $j$, Alice and Bob do the following: \begin{enumerate}
              \item \textbf{Running protocol:} In round $q$, Alice and Bob run
                the round $q$ of the protocol $\prot_\bhpc$ on the instance $(A,
                B, C, D)$ by Alice playing $P_A$, Bob playing $P_B$, and both
                Alice and Bob simulating $P_C$ and $P_D$ with no communication
                (this is possible as both Alice and Bob know $(C, D)$ entirely).
                Denote this transcript by $\Pi_q$.

           \item \textbf{Fixing pointer:} Alice and Bob compute the pointer $z_{q}$ using the fact
           that for any underlying instance $(A_x,B_x) \in (A, B) \setminus (A_{x_i},B_{x_i})$ either Alice or Bob knows the entire instance. We use $\Prot^*$ to denote the cumulative transcript generated by this step in all rounds.
           
           \item They stop if $x_i = z_{q}$.
          \end{enumerate} 
           
           
	\item The players return $\Pi_\SI := (\Pi_1,\ldots,\Pi_j,\Pi^*)$ from which $\mu(\estar)$ is calculated.
    \end{enumerate}
\end{tbox}

First, we show that the internal information cost of this protocol is small if the communication cost of $\Pi_\bhpc$ is small. This uses standard direct sum argument of information complexity. More concretely, we have the following:

\begin{claim} \label{clm:nonmdirect-sum}
    $\ICost{\pi_\SI}{\cD} = O\left(\frac{CC(\pi_\bhpc)}{m} + \frac{j \log m} m\right)$.
\end{claim}



The proof of \Cref{clm:nonmdirect-sum} is similar to that of Claim 5.6 of~\cite{AssadiCK19}. We defer it to the \Cref{sec:proof-direct-sum}. At this point, we also need a bound on how well $\Pi_\SI$ can ``almost'' solve Set-Intersection. To this end, we have the following guarantee.

\begin{claim} \label{clm:nonmaverage-coordinate}
$\pi_\SI$ externally $\eps$-solves Set-Intersection on $\cD$ where:
\[
  \eps \geq \E_{(\rE_j, \rPi_j)}\E_{x \sim \cU_\cX}
  \left[\Lambda(\mu(\rT_x \mid \rE_j, \rPi_j), \mu(\rT_x)) \right] - \frac{j} m.
\]
\end{claim}



\begin{proof} 
    This proof follows the proof of Claim 5.5 of \cite{AssadiCK19} almost verbatim. The only change is that we are working with positive triangular discrimination ($\Lambda$) instead of the total variation distance ($\Delta_{\sf TVD}$). However, as $\Lambda$ is bounded by $\Delta_{\sf TVD}$ (\Cref{prop:prop-lambda}), a similar calculation yields the claim.
    By \Cref{def:eps-solving}, $\eps = \E_{\rPi_\SI}[\Lambda(\mu(\rT \mid
    \rPi_\SI), \mu(\rT))]$. So we have:

    \begin{align*}
        \E_{\rPi_\SI}[\Lambda(\mu(\rT \mid \rPi_\SI), \mu(\rT))]
        &= \E_{\rE_j, \rPi_j, \rPi^\ast, \rI}
        [\Lambda(\mu(\rT_{x_{\rI}} \mid \rPi_\SI), \mu(\rT_{x_{\rI}}))]
        \tag{as $\rT = \rT_{x_{\rI}}$}\\
        &\hspace{-1.9cm}=
        \E_{\rE_j, \rPi_j} \E_{\rI} \E_{\rPi^\ast\mid \rE_j, \rPi_j, \rI}
        [\Lambda(\mu(\rT_{x_{\rI}} \mid \rPi_\SI), \mu(\rT_{x_i}))]
        \tag{as $\rI$ is independent of $(\rE_j, \rPi_j)$}\\
        &\hspace{-1.9cm}=
        \E_{\rE_j, \rPi_j}
        \left[\sum_i \frac 1 m \E_{\rPi^\ast\mid \rE_j, \rPi_j, \rI = i}
        [\Lambda(\mu(\rT_{x_i} \mid \rPi_\SI), \mu(\rT_{x_i}))]\right]
        \tag{as $\rI \sim \cU_{[m]}$}\\
        &\hspace{-1.9cm}=
        \E_{\rE_j, \rPi_j} \left[\sum_{x_i \in \rZ^{<j}}
          \frac 1 m \Lambda(\mu(\rT_{x_i} \mid \rPi_\SI), \mu(\rT_{x_i}))
        + \sum_{x_i \notin \rZ^{<j}} \frac 1 m \Lambda(\mu(\rT_{x_i} \mid \rE_j,
      \rPi_j), \mu(\rT_{x_i}))\right]\\
        \tag{If $x_i \notin \rZ^{< j}$ then $\rPi_\SI = (\rE_j, \rPi_j)$}\\
         &\hspace{-1.9cm}\geq
         \E_{\rE_j, \rPi_j} \left[\sum_{x_i \notin \rZ^{<j}} \frac 1 m
         \Lambda(\mu(\rT_{x_i} \mid \rE_j, \rPi_j), \mu(\rT_{x_i}))\right]
         \tag{as the first sum term is at least 0}\\
         &\hspace{-1.9cm}\geq \E_{\rE_j, \rPi_j} \E_i\left[
         \Lambda(\mu(\rT_{x_i} \mid \rE_j, \rPi_j), \mu(\rT_{x_i}))\right]
         -\frac{j}{m}. \tag{as $\Lambda$ is upper bounded by 1}
    \end{align*}
\end{proof}



\noindent
Using \Cref{clm:nonmdirect-sum} and the lower bound on $\ICost{\pi_\SI}{\cD}$ (\Cref{cor:si-ic-lb-external}), we have:
\[
\eps = O\left(\frac{CC(\pi_\bhpc)}{m^2} + \frac{j \log m} {m^2}\right)
\]
Now, using \Cref{clm:nonmaverage-coordinate}, we get:
\[
\E_{(\rE_j, \rPi_j)}\E_{x \sim \cU_\cX} \left[\Lambda(\mu(\rT_x \mid \rE_j,
\rPi_j), \mu(\rT_x)) \right] - \frac j n = O\left(\frac{CC(\pi_\bhpc)}{m^2} + \frac{j \log m} {m^2}\right),
\]
which concludes the proof of \Cref{lem:nonmprogress}.


\subsection{Putting it All Together}\label{sec:put-together}


 Given \Cref{lem:mhpc-dist}, we prove \Cref{thm:lb-mhpc} in the following way:
 Consider a protocol $\pi$ that is solving $\bmhpc_{m,r}$ with communication
 cost $o(m^2)$.
 By \Cref{lem:mhpc-dist}, we have:

 \begin{align*}
     \E_{(\rE_r, \rPi_r)}\left[\Lambda(\mu(\rZ_r \mid \rE_r, \rPi_r), \mu(\rZ_r))\right]
     &\leq 6 c \cdot \frac{\mi{\rPi}{\rA,\rC \mid \rB,\rD} + \mi{\rPi}{\rB,\rD
\mid \rA,\rC}}{m^2}\\ &\underset{(\ast)}{\leq} 12c \cdot \frac{\myCC(\pi_\bhpc)}{m^2} = o(1),
 \end{align*}

where inequality $(\ast)$ follows from \Cref{part:uniform,part:info-atmost-rv}
in \Cref{fact:it-facts}.
However, next, we show that, for a correct protocol for $\bmhpc_{m,r}$, this
quantity should be $\Omega(1)$, which contradicts the assumption that
$\myCC(\pi_\bhpc) = o(m^2)$.
An entirely similar argument derives \Cref{thm:lb-hpc} from \Cref{lem:nonmhpc-dist}.
The only difference between these two arguments is the bound on $r$: We have to
assume $r = o(\sqrt{m})$ because of the additive $r^2/m$ factor in
\Cref{lem:nonmhpc-dist} (when applied with on the $r$-th round of the protocol).




\begin{lemma} \label{lem:hpc-correctness}
    Any protocol that solves $\bmhpc_{m,r}$ on distribution $\cD_\bmhpc$, or $\bhpc_{m,r}$, on distribution $\cD_\bhpc$, with constant probability must have:
    \[
 \E_{(\rE_r, \rPi_r)}\left[\Lambda(\mu(\rZ_r \mid \rE_r, \rPi_r), \mu(\rZ_r))\right] = \Omega(1).
\]
\end{lemma}

\begin{proof}
    
 Let $b$ denote the modulus function such that $b(z_k) = i \mod 2$ iff $z_r = x_i$ or $y_i$ depending on where $r$ is even or odd. If the protocol $\pi$ computes $\bmhpc_{m,r}$ or $\bhpc_{m,r}$ with probability $2/3$, we have the following inequality:

 \begin{align*}
     2/3 &= \E_{\rE_r,\rPi_r} \Pr_{\rZ_r \mid \rE_r,\rPi_r} [b(\rZ_r) =
     \cO_\pi(\rE_r,\rPi_r)]\\
     &\underset{(1)}\leq \E_{\rE_r,\rPi_r}\left[\Pr_{\rZ_r \sim
     \cU_{\cX}}[b(\rZ_r) = \cO_\pi(\rE_r,\rPi_r)] + \|\mu(b(\rZ_r) \mid \rE_r,\rPi_r) - \cU_{0,1}\|_1 \right]\\
     &\underset{(2)}\leq \E_{\rE_r,\rPi_r}\left[\Pr_{\rZ_r \sim
     \cU_{\cX}}[b(\rZ_r) = \cO_\pi(\rE_r,\rPi_r)] + \|\mu(\rZ_r \mid \rE_r,\rPi_r) - \cU_\cX\|_1 \right]\\
     &= \frac 1 2 + \E_{\rE_r,\rPi_r}\left[\|\mu(\rZ_r \mid \rE_r,\Pi_r) - \cU_\cX\|_1 \right]
 \end{align*}
 where (1) holds because of the simple fact that $\Pr_p[\cE] \leq \Pr_q[\cE] + \|p-q\|_1$ where $p$ and $q$ are two distributions over a sample space of which $\cE$ is an event (see Footnote \ref{fn:tvd}), and (2) holds because of data processing inequality. This implies:
 \begin{align*}
     \E_{\rE_r,\rPi_r}\left[\|\mu(\rZ_r \mid \rE_r,\rPi_r) - \cU_{0,1}\|_1
     \right] \geq 1/6 &\Rightarrow \left(\E_{\rE_r,\rPi_r}\left[\|\mu(\rZ_r \mid
     \rE_r,\rPi_r) - \cU_{0,1}\|_1 \right] \right)^2\geq 1/36\\
     &\Rightarrow \E_{\rE_r,\Pi_r}\left[\|\mu(\rZ_r \mid \rE_r,\rPi_r) -
     \cU_{0,1}\|_1^2 \right] \geq 1/36\\
     &\Rightarrow \E_{\rE_r,\Pi_r}
     \left[\Lambda(\mu(\rZ_r \mid \rE_r,\Pi_r), \cU_{0,1}) \right]
     \geq 1/288 = \Omega(1),
 \end{align*}
 which completes the proof.
\end{proof}



